\documentclass{ctexart}

\usepackage{geometry}

\geometry{a4paper, right=3cm, left=3cm, top=2cm, bottom=2cm}

\usepackage[british]{babel}
\usepackage[utf8]{inputenc}
\usepackage{epstopdf}
\usepackage{csquotes}
\usepackage[hidelinks]{hyperref}
\usepackage[
    style=apa,
    backend=biber,
    sortcites=true,
    sorting=nyt,
%    isbn=false,
%    url=false,
%    doi=false,
%    eprint=false,
    hyperref=false,
    backref=false,
%    firstinits=false,
]{biblatex}

% maps apacite commands to biblatex commands
\let \citeNP \cite
\let \citeA \textcite
\let \cite \parencite

\bibliography{reference}

\usepackage[T1]{fontenc}

% Required for inserting images
\usepackage{graphicx} 

\pagestyle{plain} % 页码放下面
\title{Semantics and Syntax of Gradable Adjectives in English and Mandarin}
\date{\vspace{-10ex}}

\usepackage{enumitem}
\setlist[enumerate,1]{label=(\arabic*)}
\setlist[enumerate,2]{label=(\alph*)}

%%%%%%%%%%%%%%%%%%%%%%%%%%%%%%%%%%%%%%%%%%%%%%%%%%%

\begin{document}

\thispagestyle{empty} % 这一页清空

\begin{center}

Name: $<<$* Zhang Yi Feng *$>>$ \\
Student ID: $<<$* 1155166448 *$>>$

\end{center}

{\let\newpage\relax\maketitle}

\section{Introduction}

\setcounter{page}{1}

\noindent
Gradability is an essential concept in the studies of adjectives, which classifies adjectives into two major classes, non-gradable adjectives, like \textit{British}, etc. and gradable adjectives, like \textit{tall}, \textit{long}, etc. To grasp this distinction and particularly to account for the core characteristic of gradable adjectives, linguists conduct a great many investigations around DEGREE. Among various aspects for looking into the problem, this paper adopts the major approaches in the field of formal semantics as well as in generative grammar. The paper attempts to integrate formal semantics and generative syntax to reveal the core meaning of gradable adjectives and the syntactic constructions where they appear. The explanation for the noticeable differences between Mandarin comparisons and English comparisons will also be given.

\subsection{Starting point: Cresswell}

\noindent
Compositionality is the most important principle in formal semantics. Under the guidance of the Compositionality Principle, formal semantics focuses on determining the semantic type of constituents and figuring out the means of how they combine with each other, which is referred to as the type-driven computation \cite{pan1998}. In formal semantics, questions concerning generalised quantifiers and $\lambda$-conversion are of prominent importance and they are also largely encountered in discussions of gradable adjectives. As to degree semantics, degree is encoded in the meaning of gradable adjectives an account commonly admitted by scholars, which is initially proposed by \cite{cresswell1976}, giving rise to a third primitive semantic type $d$ besides the two primitive semantic types $e$ and $t$ in classical semantics Degree semanticists also identify three major parts of gradable adjectives: (a) a measure function $G$, mapping the target $x$ onto the abstract dimension for measurement characteristic of the gradable adjective \cite{bartsch1974}; (b) the total ordering relation $\geq$, which makes the set of scales corresponding to the abstract dimension ordered in pairs of a same direction; (c) the degree variable, indicating the value of $G(x)$. A gradable adjective of predictive use is analysed as a two-place predicate with the individual and the degree as its arguments.

The purpose of generative grammar is using definite and concise rules to generate indefinite and complex structures via formal expressions. Among a great many technical details of generative grammar, the development of the functional phrase (FP) can be taken as reference to account for the structure of degree constructions in particular. Besides the derivation of Logical Form (LF) is a cross-cutting question between syntax and semantics, which is crucial for discussions of the scope ambiguity of generalised quantifiers, almost unavoidable in debates around degree constructions.

Since Cresswell, the semantic type of gradable adjectives is largely debated, among which three major perspectives are listed here: 
(a) $<d,<e,t>>$ \cite{von1984a}, which is the so-called standard analysis; 
(b) $<e,t>$ \cite{kennedy1997}, which is referred to as the simplified analysis;
(c) $<e,<e,t>>$ \cite{rett2008}, which is referred to as the order-changed analysis.

Another doubtful topic in the field of degree constructions is the status of “more/er” degree morpheme. The major discrepancy exists in whether to regard the degree morpheme as a generalised quantifier \cite{von1984a,heim1985,bhatt2004} or to regard the degree morpheme as a functional morpheme \cite{bierwisch1989,corver1990,corver1993,corver1997a,kennedy1997} . The former including the so-called standard analysis of von Stechow, advocates that the degree morpheme lands at an adjunct position within the adjective phrase (AP), undergoing quantifier raising (QR), while the latter assumes that the degree morpheme occupies the head of the degree phrase (DegP) located above the gradable adjective.

\subsection{Standard analysis: von Stechow}

\noindent
von Stechow proposes a comprehensive constructive analysis of comparisons \cite{von1984a}, which becomes the so-called standard analysis later\cite{bale2011}. A brief account of major contents given in his analysis will be given in the following paragraphs, organized in an order to unfold semantic components of comparative constructions, where readers will see that the framework of generative grammar is compatible with this semantic analysis.

The Meaning of gradable adjectives are interpreted as a measure function plus an ordering relation. The measure function maps the individual to the dimension denoted by gradable adjectives and the ordering relation ensures that the scale of the individual exceeds the degree to compare. In generative grammar, a gradable adjective is the head of AP. It functions as a two-place predicate, with an individual typed $e$ and a degree typed $d$ as its two arguments. What deserves a note is it is von Stechow who in first seriously regards degrees denoted by symbol $d$ as one of the primitive semantic types and it is degree $d$ that captures the difference between gradable adjectives and non-gradable adjectives. Here the semantic type of gradable adjectives is manifested as $<d,<e,t>>$. Giving tall as an instance, the lexical entry of a gradable adjective is shown in \ref{ref1}, in which height denotes the measure function encoded by \textit{tall} and $\geq$ denotes the ordering relation between the individual and the degree:

\begin{enumerate}
\item \label{ref1} $[\![tall]\!] = \lambda d \lambda x. [height(x) \geq d]$
\end{enumerate}

The analysis of DegP headed by a degree morpheme is much more complicated. The semantic type of DegP is $<d,<e,t>>$. Based on the approach in generative grammar, DegP lands at an adjunct position of AP in the deep structure and then undergoes quantifier raising (QR) in the logical form (LF) to a node above the original IP inside which DegP is initially located, with the trace left denoting a $d$ type argument. The motivation for this movement is that, according to von Stechow’s analysis DegP can be regarded as isomorphic to a generalised quantifier phrase (QNP), an account under large debate afterwards.

The head of DegP, giving the comparative morpheme \textit{er} as the typical one, takes a “than phrase” (thanP) as its complement and a differential phrase (DiffP) as its specifier. Depending on various categories of complement in thanP as well as the overt appearance, such as \textit{6 feet}, or covert appearance of DiffP, \textit{er} has kinds of semantic variants:

\begin{itemize}
\item[1.]
thanP can take three types of complement:

\begin{itemize}
\item[i.]
A direct degree expression, such as than \textit{6 feet}, which denotes a degree argument typed $d$.
\item[ii.]
A comparative clause, such as \textit{than Mary is tall}, which denotes a property of degree argument, typed $<d,t>$ because according to the view of Chomsky, the comparative clause \textit{than Mary is tall} owns a deep structure looking like \textit{than} $how_i$ \textit{Mary is} $t_i$ \textit{tall} which undergoes \textit{wh}-movement, leaving the trace $t_i$ denoting a degree variable bound by $\lambda$-operator\cite{chomsky1977}.
\item[iii.]
A bare NP, such as \textit{than Mary}, regarded as a deletion from the full comparative clause \textit{than Mary is tall}, also denotes a property of degree argument typed $<d,t>$.
\end{itemize}

\item[2.]
DiffP with overt appearance denotes a degree argument, typed $d$.

\end{itemize}

Here some examples of possible semantic variants of \textit{er} are illustrated:

\begin{enumerate}[resume]
\item

\begin{enumerate}
\item John is taller than 6 feet.
\item John is taller than Mary.
\item John is taller than Mary is.
\item John is 6 inches taller than Mary is.
\end{enumerate}

\end{enumerate}

thanP takes a direct degree expression as complement in (2a), a bare NP in (2b), a comparative clause in (2c—2d). Optional DiffP only owns overt appearance in (2d). Based on the analysis above, sematic type of \textit{er} in (2a) is $<d,<<d,t>,t>>$, in (2b) as well as (2c) is $<<d,t>,<<d,t>,t>>$, in (2c) is $<<d,t>,<d,<<d,t>,t>>>$.

According to type-driven computation in formal semantics, the gradable adjective typed $<d,<e,t>>$ first combines the trace of the semantic type $d$ which is left by DegP in the process of QR, then combines the subject in the matrix clause of the semantic type $e$, outputting a $t$ type proposition with a free degree variable. $\lambda$-abstraction turns this $t$ type open proposition into a property of degree typed $<d,t>$ which saturates DegP typed $<<d,t>,t>$. Finally, a $t$ type proposition is made out. Here we have successfully recapitulated the procedure in von Stechow’s analysis of how to derive a comparative construction.

\subsection{Simplified version: Kennedy}

The research of Kennedy investigates what the so-called gradable adjectives refer to as well as in what kind of constructions they appear, which is referred to as degree constructions in his paper\cite{kennedy1997}. In this research, Kennedy combines semantics and syntax to reveal the essence of this largely discussed topic.

As to the semantic interpretation of gradable adjectives, there are two major schools: (a) scalar analyses, in which the set of scales corresponding to a gradable adjective is compatible with the partial ordering relation\cite{cresswell1976,bierwisch1989}; (b) vague predicate analyses, in which the extensions of gradable adjectives need to be determined according to the context\cite{ginet1973,kamp2013,klein1980}. Kennedy’s analysis falls into the general family of scalar analyses.

However, Kennedy also makes some revolution to the classical scalar analysis, particularly as to the semantic type of gradable adjectives, as well as the interpretation of degrees.

Firstly, classical scalar analyses characterize the core meaning of gradable adjectives as an ordering relation between the target and the degree. Most scholars employing scalar approaches also believes the free variable degree to be bound by a quantificational expression\cite{hellan1981,hoeksema1983,von1984a}. In contrast, Kennedy assumes that gradable adjectives denote measure function, giving rise to his non-quantificational analysis of degree constructions.

Kennedy identifies three semantic constituents in degree constructions\cite{russell1905}: (a) a reference value $G(x)$, expressing the projection from the target $x$ to the scale associated with the gradable adjective $G$; (b) a standard value $s$, which is indicated via measure phrasess or the comparative clause landing at the complement position of the preposition “than”; (c) degree relation $R$, a partial ordering relation encoded by degree morphemes, elements of \textit{\{er/more, less, as, too, enough, so, how, ...\}}, which compares the reference value and the standard value. The formalization of degree constructions is shown in \ref{ref3}:

\begin{enumerate}[resume]
\item \label{ref3} $\lambda x.[R(G(x))(s)]$
\end{enumerate}

In terms of the syntactic structure of degree constructions, Kennedy projects the degree phrase (DegP) headed by a degree morpheme as a functional phrase (FP) above the adjective phrase (AP) headed by a gradable morpheme, an account which is developed from approaches where adjectives are projected to an extended functional structure\cite{abney1987,corver1990,corver1997b,grimshaw1991}. The syntactic tree structure of degree constructions is shown in \ref{ref4}:

\begin{enumerate}[resume]
\item \label{ref4}
\end{enumerate}

\begin{figure}[h]
\centering
\includegraphics[width=0.5\textwidth]{pic/pic1.png}
\end{figure}

Typically, scholars think that the meaning of gradable adjectives is composed of two major parts: a measure function and an ordering relation\cite{carston2013,horn1992}.

In contrast, Kennedy’s analysis excludes the relational component from the meaning of gradable adjectives, rather assigning this semantic part to the degree morphology. This division of labour between the gradable adjective and the degree morphology is also reflected by the AP/DegP extended projection. Therefore, Kennedy’s analysis revises the semantic type of gradable adjectives from the traditional version $<d,<e,t>>$ to the simplified version $<e,t>$, making the complexity transmitted to the degree morphology.

Secondly, Kennedy makes refinement to the interpretation of degrees. He argues that degrees should be formalized as extents\cite{seuren1978,von1984b,bierwisch1989,lobner2010} rather than as points on a scale, which is different from classical scalar analyses. Based on this extent-scalar view, Kennedy assumes that gradable adjectives characteristic of adjectival polarity can be divided into positive adjectives, such as \textit{tall}, and negative ones, such as \textit{short}. According to this distinction, the anomaly of comparisons made between antonymous adjectives as shown in \ref{ref5} can be explained by the undefined ordering relation for extents of opposite polarity:

\begin{enumerate}[resume]
\item \label{ref5}
*Carmen is taller than Mike is short.
\end{enumerate}

%%%%%%%%%%%%%%%%%%%%%%%%%%%%%%%%%%%%%%%%%%%%%%%%%%%
%%%%%%%%%%%%         Ch2         %%%%%%%%%%%%%%%%%%
%%%%%%%%%%%%%%%%%%%%%%%%%%%%%%%%%%%%%%%%%%%%%%%%%%%


\section{DegP hypothesis and Deg-Shell}

\subsection{DegP hypothesis}

% 刚出现DegP的时代，在adjP的adjunt位置上，需要经历Quantifier raising，比较支持这个的有crosswell，von那个sb，后来，从abny开始，提出了degP是一个functional projection（相当于强化了句法功能），并且是located above AP，

% 以前的 degP是aj
\begin{enumerate}[resume]
\item
$[_{AP} [_{A^{\prime}} DegP[_{A^{\prime}} A^0]]]$
\end{enumerate}

% 现在degP是
\begin{enumerate}[resume]
\item
$[_{DegP} [_{Deg^{\prime}} Deg^0 \enspace AP]]$
\end{enumerate}

% Abney根据IP和VP的关系，抄袭了一个DP和NP的关系（DP在NP之上），顺便抄袭了一下DegP和AP的关系，指DegP在AP之上。
% gerund 

% 句法上，DegP head Gradable adj -> lexcial projection， DegP是funcitonal projection
% 语义上，adj有个d需要人填

To study the linguistic features of gradable adjectives, Abney established the DegreeP Hypothesis in his research\cite{abney1987}. Before his work, degree words are always treated as occupying the specifier position of AP\cite{bowers1975,jackendoff1977}. In DegreeP Hypothesis, the comparative and superlative forms of adjectives are analyzed as a functional category called Degree Phrase(DegP), taking the lexical projection AP as its complement.

\subsection{DegP-Shell}

% Larson 1988 VP-shell 1991？ DegP shell

\subsection{DegP Analysis in Sepcific Western Language}

There is lots of work about DegP analysis on specific languages. In most cases, language-specific analyses always apply language-specific rules, which makes their conclusions hard to be extended to different languages such as Mandarin, but the approaches they apply to problems are still valuable. So here I give two examples to show how previous researchers conduct language-specific analyses based on the specific languages which they are familiar with.

In Corver's research in 1997, he investigates the phrase structure of the adjective projection in Dutch language \cite{corver1997a} based on DegP hypothesis. He found empirical evidence about distinctions between DegP and QP(means quantifier-like degree phrase), \ref{degPAndQp}\ref{degPAndQpDegP} shows examples of DegP and \ref{degPAndQp}\ref{degPAndQpQp} gives examples of QP. This distinction also leads different phrasal structure, which is shown in\ref{ref7}. 

\begin{enumerate}[resume]
\item \label{degPAndQp}
\begin{enumerate}
\item \label{degPAndQpDegP} 
zo, te, hoe, even \\
zo, too, how, as
\item \label{degPAndQpQp}
meer, minder, genoeg \\
more, less, enough
\end{enumerate}
\end{enumerate}

\begin{enumerate}[resume]
\item \label{ref7}

\begin{enumerate}
\item \label{ref7a} 
$[_{DegP} \enspace te[_{QP} \enspace e[_{AP}[_{A^{\prime}} \enspace lang]]]]$ \\
$[_{DegP} \enspace too[_{QP} \enspace e[_{AP}[_{A^{\prime}} \enspace tall]]]]$
\item \label{ref7b} 
$[_{QP} \enspace minder[_{AP}[_{A^{\prime}}\enspace lang]]]$ \\
$[_{QP} \enspace less[_{AP}[_{A^{\prime}}\enspace tall]]]$

\end{enumerate} 
    
\end{enumerate} 

In \ref{ref7}\ref{ref7a}, syntactic structure shows that the DegP head selects a QP and QP's head in turn selects AP. \ref{ref7}\ref{ref7b} shows that adjective phrases are introuduced by quantifier-like elements and the DegP is missing. The author tried to merge two structures into a single elegant structure. To make this happen, some language-specific evidence must be discovered to support the so-called "$A^0$ to $Q^0$ raising", which raises the adjectival head into the empty functional Q-position. In Dutch, there are serval phenomena suggesting the existence of this movement. Take one example as \textit{genoeg}-inversion, which is illustrated in \ref{genoeg}\ref{genoeg_b}, and the structure of \textit{genoeg}-inversion is shown in \ref{genoeg}\ref{genoeg_c}, which overtly shows the existence of $A^0$ to $Q^0$ raising in Dutch.

\begin{enumerate}[resume]
\item \label{genoeg}
\begin{enumerate}

\item 
*Jan is [genoeg bang daarvoor] \\
Jan is enough afraid there-of \\
Jan is afraid enough of it

\item \label{genoeg_b} 
Jan is [bang genoeg daarvoor] \\
Jan is afraid enough there-of \\
Jan is afraid enough of it

\item \label{genoeg_c} $[_{QP} \enspace bang_i + genoeg[_{AP} \enspace t_i \enspace daarvoor]]$

\end{enumerate} 
\end{enumerate} 

Besides, Corver discussed the way to build a more complete syntactic structure. By summarizing serval language-specific rules, he found that there has to be a AgrP in the extended adjectival projection. For one side, if the external argument is assumed to originate in the specifier of AP, there should be a syntactic position so that the sbject can be moved in; and for other side, AgrP gives a better solution to explain some Dutch-specific movement operations. After all his work, Corver's final articulated structure whth AgrP as bellow:

\begin{enumerate}[resume]
\item

$[_{AgrP}[_{Agr^{\prime}}[_{DegP}[_{Deg^{\prime}}Deg[_{QP}[_{Q^{\prime}} Q [_{AP}DP[_{X^{\prime}}A \enspace XP]]]]]]Agr]]$
    
\end{enumerate}

Another example of language-specific analysis is Svenonius and Kennedy's studied about Null Degree Questions in Scandinavian \cite{svenonius2006}. The Null Degree Question is such a phenomenon that, in some dialects in Scandinavian, a degree sentence can only have degree operator but no \textit{wh}-word. For example, \ref{refNDQ}\ref{refNDQa} is a sentence in Icelandic, and \ref{refNDQ}\ref{refNDQb} shows the mapping from Icelandic words to English words. What should be noticed is that the Icelandic word \textit{hversu} is mapped to two English words \textit{how much}, but actually \textit{hversu} only used as a degree operator, but not a manner adverbial. This phenomenon makes sentence '\textit{how old are you?}' sound like '\textit{old are you?}'. 

\begin{enumerate}[resume]
\item \label{refNDQ}
\begin{enumerate}
\item \label{refNDQa} Hversu gammall ertu? 
\item \label{refNDQb} \textit{how.much old are.you}
\item \label{refNDQc} How old are you?
\end{enumerate}
\end{enumerate}

Besides the expression above, there is another kind of expression in Icelandic, by which speaker can front a question word \textit{Hva\dh}, and leave the adjective in original poisiton, shown in \ref{refHva}. To explain this phenomenon, authers make a postulation of null operator in Null Degree Question, whose structure shows in \ref{refOp1}.

\begin{enumerate}[resume]
\item \label{refHva}
\begin{enumerate}
\item Hva\dh$\enspace$ertu gammall? 
\item \textit{what are.you old}
\item How old are you?
\end{enumerate}
\end{enumerate}

\begin{enumerate}[resume]
\item \label{refOp1}
\begin{enumerate}
\item \label{refOp1a} $[_{CP}Op_1 \enspace er_2[_{IP} du_3 \enspace t_2 [_{VP} t_2[_{AP} t_3 \enspace t_1 \enspace gammel]]]]$ \\
\textit{Op are you old} \\
How old are you?
\item \label{refOp1b} $[_{CP}$Hva\dh$\enspace er_2[_{IP} -tu_2 \enspace t_2 [_{VP} t_2[_{AP} t_3 \enspace t_1 \enspace gammel]]]]$ \\
\textit{What are you old} \\
How old are you?
\end{enumerate}
\end{enumerate}

Compare \ref{refOp1}\ref{refOp1a} and \ref{refOp1}\ref{refOp1b}, it is obvious to see the fact that, the operator $Op_1$ and the question word \textit{Hva\dh} \enspace have the same behaviour as \textit{wh}-movement, and they both follow the rule of locality conditions. That is the evidence makes authors belive the existence of null operator, which is $Op_1$ in \ref{refOp1}\ref{refOp1a}.

% (52)

\subsection{DegP Analysis in Mandarin}

%%%%%% Mandarin transitive comparatives and the grammar of measurement (2012)
%%%%%% Thomas Grano & Chris Kenndy

% %%%%%%%%%%%%%%%%%%%%%%%%%%%%%%%%%%%%%%%%%%%%%%%%%%%%%%%%%%%%%%%%%%%%%%
% %%%%%%%%%%%%%%%%%%%%%%%%%%%% Introduction %%%%%%%%%%%%%%%%%%%%%%%%%%%%
% %%%%%%%%%%%%%%%%%%%%%%%%%%%%%%%%%%%%%%%%%%%%%%%%%%%%%%%%%%%%%%%%%%%%%%

% canonical comparative in Mandarin -> morpheme "bi"
% A bi B gao (yi dian)

% transitive comparative -> "bi" absent, gradable adj as gradable predicate
% A gao B yi dian

% transitive compartive is not always abailable
% * A piao liang B yi dian
% Xiang(2005) gives two reason:
% - transitive comparative must contain an overt differential term
% - transitive comparative's gradable predicate must uses "conventional-measurement-system-defined" scale.
% ("conventional measurement system defined" means "gao san gong fen" is ok, "piao liang san du" is ill form).
% The background reason of weather a predicate can be used in transitive comparative is weather the comparative allows a differential measure phrase.

% Based on theories above, Xiang depart gradable predicate into two parts:
% - Grammatical in Transitive Comparative: gao, ai, pang, show, qing, chang, duan,..
% - Ungrammatical in Transitive Comparative: piao liang, you qu, gan jing, shu fu,..

% (Doetjes 1997, Neeleman 2004, Kennedy & McNaaly 2005b) Vague degree terms like "a bit, a little" can function either as true measure phrases or as a degree modifiers, which have different syntactic and semantic properties.
% a bit heavy
% * 40 kg heavy (disallow measure phrase)
% very heavy (allow degree modifier very)
% deepen the hole a bit
% deepen the hole 2 feet (allow measure phrase)
% * deepen the hole very (disallow degree modifer)

% (full discussion of difference between measure phrase & degree modifies is beyond this paper -> adopt the result of Kennedy & McNally 2005b -> measure phrase saturate a degree argument position while modifiers map properties to properties.)

% main purpose of this paper:
% - explain why a differential phrase is obligatory in Mandarin transitive comparatives (Noted by Xiang).
% - explain why the adjective must use a scale with a defined measurement system in Mandarin transitive comparatives (Noted by Xiang).
% - argue that there are two case assigners for standards of comparsion in Mandarin.
% -- standard marker "bi"
% -- covert morpheme like null operator in dialect in Norwegian, which is involved in the projection of measure phrase.

% %%%%%%%%%%%%%%%%%%%%%%%%%%%%%%%%%%%%%%%%%%%%%%%%%%%%%%%%%%%%%%%%%%%%%%
% %%%%%% The synatax and semantics of gradability and measurement %%%%%%
% %%%%%%%%%%%%%%%%%%%%%%%%%%%%%%%%%%%%%%%%%%%%%%%%%%%%%%%%%%%%%%%%%%%%%%

% Rule(P227): If a measure phrase denotes a degree on a scale that is distinct form the one used by a particular adjective, then composiiton will be ruled out as a kind of selectional restriction violation.
% * 20 kg gao (gao do not accept kg)
% * 5 du piao liang (piao liang use scales without defined measurement systems)

% Rule(P227): Measure phrase do not combine with negative-pole adjectives.
% * 1 mi duan
% 1 mi chang

% Problem: why "1 meter short" is ill form but "1 meter long" is correct. Futhermore, why "1 meter shorter" is correct.
% Old schools belives that it is beacuse those expression do not include a fixed origin point. (Seuren 1978, von Stechow 1984b, Bierwisch 1989, Kennedy 2001, Sassoon 2010).

% A solution from Schwarzschild 2005: change measure phrase's analysis from "a point in number axis" to "a gap or an interval in number axis".
% Whcih means, this change lexical entry of mearsure phrase, forn "d" to "sets of d". thus, lexcial adjectives are type <d, <e, t>> and can not accept mearsure phrase as input.
% Which means, "2 cm heavy" is ill-formed, but "2 meter high" is also ill-formed. Schwarzschild belive this is a good result because he thinks "measure phrase + postive adjective" is an idiosyncratic and language-specific.
% So how to resolove the problem that measure phrase can not be a input to adjectives? See (25) Homonym Rule.

% Another solution form Svenonius & Kenndy 2006: analyze gradable adjectives as measure function: expression of type <e, d>.
% Which leads a problem immediately, if measure phrases are type d, they can not combine with gradable adjectives, this is ruled out on type-theoretic gounds.
% Svenonius & Kenndy belive that this type mismatch can be resolved through the mediation of functionl morphplogy. So they propose that, among the inventory of degree morphemes is null head "u", whose semantic and syntactic function is to introduce a degree arguement.
% [u] = \lambda g \lambda d \lambda x. g(x) \geq d
% [tall] = height
% [u]([tall]) = \lambda d \lambda x. height(x) \geq d
% u combines freely with comparative adjectives, but only idiosyncratically with non-comparative adjectives, which is language-specific.
% which means, "10 kg heavier" is ok, "10 meter longer" is ok. "10 ky heavy" is not ok because English do not allow it, "10 meter long" is ok. In German they both ok, in Japanese neither is ok.
% u's position: DegP head -> [_{DegP} Mp [_{DegP^{\prime}} u \enspace AP]] -> 2 cm(MP) u(u) tall(AP).

% Semantically, authors follow Faller(2000), Kenndey & McNally(2005a), Svenonius & Kenndy(2006), Kennedy & Levin(2008), who analyze comparative adjectives as special kinds of measure function called "difference function".
% heavy: mapping individuals to values on weight scale.
% heavier than Kim: mapping individuals to the subpart of the weight scale (> Kim's weight).
% The comparative morphology turns measure function of type <e,d> into a difference fucntion <e,<e,d>>.
% [comp] = \lambda g_<e,d> \lambda y. (g_y)^{\uparrow}
% [heavier_{comp}]([than Kim]) = (weight_{Kim})^{\uparrow}

% %%%%%%%%%%%%%%%%%%%%%%%%%%%%%%%%%%%%%%%%%%%%%%%%%%%%%%%%%%%%%%%%%%%%%%
% %%%%%%%%%%%% Analysis of Mandarin transitive comparatives %%%%%%%%%%%%
% %%%%%%%%%%%%%%%%%%%%%%%%%%%%%%%%%%%%%%%%%%%%%%%%%%%%%%%%%%%%%%%%%%%%%%

% Mandarin comparative need morpheme "bi" as case assigner.
% * A gao B
% A bi B gao

% So why Case violation can be repaired not just only by morpheme "bi", but also a measure phrase.
% A gao B 2cm.
% Author belive morpheme u takes the place of Case assigner under this case.


















\section{My Theory}

\newpage

\printbibliography

\end{document}
